\usetikzlibrary{arrows.meta, fit, backgrounds, positioning, calc}
\begin{figure}[htbp]
\centering
\resizebox{\textwidth}{!}{%
\begin{tikzpicture}[
    >=Stealth,
    % 科学问题（固定尺寸）
    prob/.style={draw, rounded corners=3pt, text width=3.6cm,
        minimum height=1.2cm, align=center, font=\small, inner sep=5pt,
        thick, fill=blue!8, draw=blue!60!black},
    % 研究工作（固定尺寸）
    work/.style={draw, rounded corners=3pt, text width=8.6cm,
        minimum height=0.9cm, align=center, font=\small, inner sep=5pt,
        thick, fill=teal!8, draw=teal!60!black},
    % 箭头
    arr/.style={->, thick, draw=black!40},
    prog/.style={->, thick, draw=black!25, dashed},
    lbl/.style={font=\footnotesize, fill=white, inner sep=1.5pt, text=black!55},
]

% === 精确坐标 ===
\def\px{2.6}       % 科学问题列中心x
\def\wx{10.6}      % 研究工作列中心x
\def\rowgap{2.0}   % 行间距
\def\subgap{1.1}   % W3-W4子间距

% 行y坐标：4个问题等间距
\pgfmathsetmacro{\yA}{0}
\pgfmathsetmacro{\yB}{-\rowgap}
\pgfmathsetmacro{\yC}{-2*\rowgap}
\pgfmathsetmacro{\yD}{-3*\rowgap}

% W3/W4 分列在P3两侧
\pgfmathsetmacro{\yW}{0.5*\subgap}

% ===== 左侧竖版标题 =====
\node[rotate=90, font=\large\bfseries, text=violet!70!black]
    at (-1.0, {0.5*(\yA+\yD)}) {端到端语音到语音生成};
\node[rotate=90, font=\normalsize, text=black!45]
    at (-0.1, {0.5*(\yA+\yD)}) {目标：高质量、低延迟};

% ===== 科学问题列（等间距） =====
\node[prob] (p1) at (\px, \yA) {\textbf{科学问题一}\\预训练模型的组合与利用};
\node[prob] (p2) at (\px, \yB) {\textbf{科学问题二}\\建模能力与解码效率的矛盾};
\node[prob] (p3) at (\px, \yC) {\textbf{科学问题三}\\多模态序列的同步流式生成};
\node[prob] (p4) at (\px, \yD) {\textbf{科学问题四}\\双向语音流的统一序列建模};

% ===== 递进箭头 =====
\draw[prog] (p1.south) -- node[lbl, right=2pt] {解码加速} (p2.north);
\draw[prog] (p2.south) -- node[lbl, right=2pt] {同步生成} node[lbl, left=2pt] {大模型} (p3.north);
\draw[prog] (p3.south) -- node[lbl, right=2pt] {全双工} (p4.north);

% ===== 研究工作列 =====
\node[work] (w1) at (\wx, \yA) {基于模块化组合的端到端语音到语音翻译模型（第三章）};
\node[work] (w2) at (\wx, \yB) {基于有向无环图的语音到语音翻译模型（第四章）};
\node[work] (w3) at (\wx, {\yC+\yW}) {基于非自回归流式解码的实时语音对话大模型（第五章）};
\node[work] (w4) at (\wx, {\yC-\yW}) {基于自回归流式解码的高质量语音对话大模型（第六章）};
\node[work] (w5) at (\wx, \yD) {基于多信道交织的全双工语音对话模型（第七章）};

% ===== 连线 =====
\draw[arr] (p1.east) -- (w1.west);
\draw[arr] (p2.east) -- (w2.west);
\draw[arr] (p3.east) -- (w3.west);
\draw[arr] (p3.east) -- (w4.west);
\draw[arr] (p4.east) -- (w5.west);

\end{tikzpicture}%
}
\bicaption{\enspace 本文研究内容概览}{\enspace Overview of the research content}
\label{fig:research_overview}
\end{figure}
